% GENERATED FILE. Edit canonical lessons, then run npm run generate:books.
% canonical-source-hash: fnv1a64:48df6d0c7ccfd8ee
% canonical-lessons: KA-C337-nidhana, KA-C337-antyakriye, KA-C337-smasana, KA-C337-sraddhanjali, KA-C337-santaisu

\chapter{Saying Goodbye}
\label{ch:ka-saying-goodbye}
\hlchaptermodality{337}

\begin{chapteropening}
\textbf{By the end of this chapter:} \emph{I can say a passing away, a funeral, a cremation ground, a tribute and to comfort.}

Complete the last lesson of chapter 337: I can say a passing away, a funeral, a cremation ground, a tribute and to comfort.
\end{chapteropening}

\section[\texorpdfstring{\kn{ನಿಧನ} (nidhana)}{ನಿಧನ (nidhana)}]{\kn{ನಿಧನ} (nidhana) — a passing away; a death}

\label{lesson:KA-C337-nidhana}



\begin{quote}
Before the new one: say the Kannada for to fight, then the Kannada for to retreat.
\end{quote}

\subsection*{You'll want to know: \kn{ನಿಧನ}}
\textbf{\kn{ನಿಧನ}} — \emph{nidhana} — \textquotedblleft{}a passing away; a death\textquotedblright{}.

\begin{quote}
\kn{ಹಿರಿಯ} \kn{ಲೇಖಕರ} \kn{ನಿಧನದ} \kn{ಸುದ್ದಿ} \kn{ಬೆಳಗ್ಗೆ} \kn{ಬಂತು}.
\end{quote}

\begin{quote}
\emph{hiriya lēkhakara nidhanada suddi beḷagge bantu.}
\end{quote}

\begin{quote}
The news of the senior writer's passing came in the morning.
\end{quote}

The polite, formal word for a death, used in the news: \kn{ಅವರು} \kn{ನಿಧನರಾದರು}, he passed away.

\begin{cousinweb}
\textbf{\kn{ನಿಧನ}} (\emph{nidhana}) is borrowed from Sanskrit \emph{nidhana}, \textquotedblleft{}settling down, residence\textquotedblright{}, and also \textquotedblleft{}conclusion, death\textquotedblright{} (Monier-Williams).
\end{cousinweb}

\begin{grammarlens}[title={What you've built}]
One more word for telling what happened.
\end{grammarlens}

\subsection*{Your turn}
\begin{itemize}
\raggedright
  \item \emph{Say it:} \emph{nidhana}
  \item \emph{Say it:} \emph{nidhana}, once more
  \item \emph{Say it:} \emph{hiriya lēkhakara nidhanada suddi beḷagge bantu.} — The news of the senior writer's passing came in the morning.
  \item \emph{Recall:} say the Kannada for to fight, then the Kannada for to retreat, then say \emph{nidhana} again
\end{itemize}

\begin{tcolorbox}[breakable,colback=teal!4,colframe=teal!35!black,title={Before you move on}]
What does \kn{ನಿಧನ} mean? (\textquotedblleft{}A passing away; a death\textquotedblright{}.) Say it once more.
\end{tcolorbox}
\section[\texorpdfstring{\kn{ಅಂತ್ಯಕ್ರಿಯೆ} (antyakriye)}{ಅಂತ್ಯಕ್ರಿಯೆ (antyakriye)}]{\kn{ಅಂತ್ಯಕ್ರಿಯೆ} (antyakriye) — a funeral}

\label{lesson:KA-C337-antyakriye}



\begin{quote}
Before the new one: say the Kannada for to retreat, then the Kannada for a passing away.
\end{quote}

\subsection*{You'll want to know: \kn{ಅಂತ್ಯಕ್ರಿಯೆ}}
\textbf{\kn{ಅಂತ್ಯಕ್ರಿಯೆ}} — \emph{antyakriye} — \textquotedblleft{}a funeral\textquotedblright{}.

\begin{quote}
\kn{ಅಜ್ಜನ} \kn{ಅಂತ್ಯಕ್ರಿಯೆಗೆ} \kn{ಊರಿನ} \kn{ಜನರೆಲ್ಲ} \kn{ಬಂದರು}.
\end{quote}

\begin{quote}
\emph{ajjana antyakriyege ūrina janarella bandaru.}
\end{quote}

\begin{quote}
The whole village came to Grandfather's funeral.
\end{quote}

\kn{ಅಂತ್ಯ}, the last, and \kn{ಕ್ರಿಯೆ}, a rite: the last rites.

\begin{cousinweb}
\textbf{\kn{ಅಂತ್ಯಕ್ರಿಯೆ}} (\emph{antyakriye}) is borrowed from Sanskrit \emph{antyakriyā}, \textquotedblleft{}the last rites, a funeral\textquotedblright{}, from \emph{antya} \textquotedblleft{}last\textquotedblright{} and \emph{kriyā} \textquotedblleft{}an act, a rite\textquotedblright{} (Monier-Williams).
\end{cousinweb}

\begin{grammarlens}[title={What you've built}]
One more word for telling what happened.
\end{grammarlens}

\subsection*{Your turn}
\begin{itemize}
\raggedright
  \item \emph{Say it:} \emph{antyakriye}
  \item \emph{Say it:} \emph{antyakriye}, once more
  \item \emph{Say it:} \emph{ajjana antyakriyege ūrina janarella bandaru.} — The whole village came to Grandfather's funeral.
  \item \emph{Recall:} say the Kannada for to retreat, then the Kannada for a passing away, then say \emph{antyakriye} again
\end{itemize}

\begin{tcolorbox}[breakable,colback=teal!4,colframe=teal!35!black,title={Before you move on}]
What does \kn{ಅಂತ್ಯಕ್ರಿಯೆ} mean? (\textquotedblleft{}A funeral\textquotedblright{}.) Say it once more.
\end{tcolorbox}
\section[\texorpdfstring{\kn{ಸ್ಮಶಾನ} (smaśāna)}{ಸ್ಮಶಾನ (smaśāna)}]{\kn{ಸ್ಮಶಾನ} (smaśāna) — a cremation ground; a cemetery}

\label{lesson:KA-C337-smasana}



\begin{quote}
Before the new one: say the Kannada for a passing away, then the Kannada for a funeral.
\end{quote}

\subsection*{You'll want to know: \kn{ಸ್ಮಶಾನ}}
\textbf{\kn{ಸ್ಮಶಾನ}} — \emph{smaśāna} — \textquotedblleft{}a cremation ground; a cemetery\textquotedblright{}.

\begin{quote}
\kn{ಸ್ಮಶಾನ} \kn{ಊರಿನ} \kn{ಹೊರಗೆ} \kn{ನದಿಯ} \kn{ದಡದಲ್ಲಿ} \kn{ಇದೆ}.
\end{quote}

\begin{quote}
\emph{smaśāna ūrina horage nadiya daḍadalli ide.}
\end{quote}

\begin{quote}
The cremation ground is outside the village, on the river bank.
\end{quote}

Used for both a cremation ground and a burial ground.

\begin{cousinweb}
\textbf{\kn{ಸ್ಮಶಾನ}} (\emph{smaśāna}) is borrowed from Sanskrit \emph{śmaśāna}, \textquotedblleft{}a place for burning the dead, a burial place\textquotedblright{} (CDIAL 12658).
\end{cousinweb}

\begin{grammarlens}[title={What you've built}]
One more word for telling what happened.
\end{grammarlens}

\subsection*{Your turn}
\begin{itemize}
\raggedright
  \item \emph{Say it:} \emph{smaśāna}
  \item \emph{Say it:} \emph{smaśāna}, once more
  \item \emph{Say it:} \emph{smaśāna ūrina horage nadiya daḍadalli ide.} — The cremation ground is outside the village, on the river bank.
  \item \emph{Recall:} say the Kannada for a passing away, then the Kannada for a funeral, then say \emph{smaśāna} again
\end{itemize}

\begin{tcolorbox}[breakable,colback=teal!4,colframe=teal!35!black,title={Before you move on}]
What does \kn{ಸ್ಮಶಾನ} mean? (\textquotedblleft{}A cremation ground; a cemetery\textquotedblright{}.) Say it once more.
\end{tcolorbox}
\section[\texorpdfstring{\kn{ಶ್ರದ್ಧಾಂಜಲಿ} (śraddhāñjali)}{ಶ್ರದ್ಧಾಂಜಲಿ (śraddhāñjali)}]{\kn{ಶ್ರದ್ಧಾಂಜಲಿ} (śraddhāñjali) — a tribute (to the dead)}

\label{lesson:KA-C337-sraddhanjali}



\begin{quote}
Before the new one: say the Kannada for a funeral, then the Kannada for a cremation ground.
\end{quote}

\subsection*{You'll want to know: \kn{ಶ್ರದ್ಧಾಂಜಲಿ}}
\textbf{\kn{ಶ್ರದ್ಧಾಂಜಲಿ}} — \emph{śraddhāñjali} — \textquotedblleft{}a tribute (to the dead)\textquotedblright{}.

\begin{quote}
\kn{ಶಾಲೆಯ} \kn{ಮಕ್ಕಳು} \kn{ಸೈನಿಕರಿಗೆ} \kn{ಶ್ರದ್ಧಾಂಜಲಿ} \kn{ಅರ್ಪಿಸಿದರು}.
\end{quote}

\begin{quote}
\emph{śāleya makkaḷu sainikarige śraddhāñjali arpisidaru.}
\end{quote}

\begin{quote}
The school children paid tribute to the soldiers.
\end{quote}

A tribute is offered to the dead, often with a minute of silence.

\begin{cousinweb}
\textbf{\kn{ಶ್ರದ್ಧಾಂಜಲಿ}} (\emph{śraddhāñjali}) is built from Sanskrit: \emph{śraddhā}, \textquotedblleft{}faith, trust\textquotedblright{} (CDIAL 12678), and \emph{añjali}, \textquotedblleft{}the hollowed hands placed together\textquotedblright{} (CDIAL 171), an offering made with joined hands.
\end{cousinweb}

\begin{grammarlens}[title={What you've built}]
One more word for telling what happened.
\end{grammarlens}

\subsection*{Your turn}
\begin{itemize}
\raggedright
  \item \emph{Say it:} \emph{śraddhāñjali}
  \item \emph{Say it:} \emph{śraddhāñjali}, once more
  \item \emph{Say it:} \emph{śāleya makkaḷu sainikarige śraddhāñjali arpisidaru.} — The school children paid tribute to the soldiers.
  \item \emph{Recall:} say the Kannada for a funeral, then the Kannada for a cremation ground, then say \emph{śraddhāñjali} again
\end{itemize}

\begin{tcolorbox}[breakable,colback=teal!4,colframe=teal!35!black,title={Before you move on}]
What does \kn{ಶ್ರದ್ಧಾಂಜಲಿ} mean? (\textquotedblleft{}A tribute (to the dead)\textquotedblright{}.) Say it once more.
\end{tcolorbox}
\section[\texorpdfstring{\kn{ಸಂತೈಸು} (santaisu)}{ಸಂತೈಸು (santaisu)}]{\kn{ಸಂತೈಸು} (santaisu) — to comfort; to console}

\label{lesson:KA-C337-santaisu}



\begin{quote}
Before the new one: say the Kannada for a cremation ground, then the Kannada for a tribute.
\end{quote}

\subsection*{You'll want to know: \kn{ಸಂತೈಸು}}
\textbf{\kn{ಸಂತೈಸು}} — \emph{santaisu} — \textquotedblleft{}to comfort; to console\textquotedblright{}.

\begin{quote}
\kn{ನೆರೆಯವರು} \kn{ಅಳುತ್ತಿದ್ದ} \kn{ಅಜ್ಜಿಯನ್ನು} \kn{ಸಂತೈಸಿದರು}.
\end{quote}

\begin{quote}
\emph{nereyavaru aḷuttidda ajjiyannu santaisidaru.}
\end{quote}

\begin{quote}
The neighbours consoled the grandmother, who was crying.
\end{quote}

Past (I …): \kn{ಸಂತೈಸಿದೆ}, \emph{santaiside}.

The person comforted takes -\kn{ಅನ್ನು}.

\begin{cousinweb}
\textbf{\kn{ಸಂತೈಸು}} (\emph{santaisu}) is usually traced to Sanskrit \emph{sāntva}, \textquotedblleft{}consolation, gentle words\textquotedblright{} (Monier-Williams), with the Kannada verb ending -\emph{isu}.
\end{cousinweb}

\begin{grammarlens}[title={What you've built}]
One more word for telling what happened.
\end{grammarlens}

\subsection*{Your turn}
\begin{itemize}
\raggedright
  \item \emph{Say it:} \emph{santaisu}
  \item \emph{Say it:} \emph{santaisu}, once more
  \item \emph{Say it:} \emph{nereyavaru aḷuttidda ajjiyannu santaisidaru.} — The neighbours consoled the grandmother, who was crying.
  \item \emph{Recall:} say the Kannada for a cremation ground, then the Kannada for a tribute, then say \emph{santaisu} again
\end{itemize}

\begin{tcolorbox}[breakable,colback=teal!4,colframe=teal!35!black,title={Before you move on}]
What does \kn{ಸಂತೈಸು} mean? (\textquotedblleft{}To comfort; to console\textquotedblright{}.) Say it once more.
\end{tcolorbox}
